\documentclass[12pt]{article}

\usepackage{ucs}     % unicode 
\usepackage[utf8x]{inputenc}  % utf-8
\usepackage[ngerman]{babel}  % new german spelling 
\usepackage{graphicx}   % use graphics 
\usepackage{fancyhdr}   % header and footer
\usepackage{setspace}
\usepackage{url}
\usepackage{tikz}
\usepackage{tikz-qtree}
\usepackage[printonlyused]{acronym}
\usepackage{a4wide}
\usepackage{natbib}	% bibstyle 
\usepackage[section]{placeins}	% \FloatBarrier


\setcounter{tocdepth}{2}
\begin{document}

\thispagestyle{empty} 

\begin{center}
	\topskip0pt
	\vspace*{\fill}
	
	\Huge{\textbf{Seminararbeit}}\\
	\vspace{1.5cm}
	
	\Large{\textbf{Automatische Bestimmung semantischer Rollen im deutschen
	Korpus SALSA 2.0}}\\
	\vspace{1cm}
		
	\includegraphics[scale=0.3]{images/logo_hu.png}
	\vspace{1cm}

	\begin{Large}
		Institut für Informatik \& Institut für Linguistik\\
		Humboldt-Universität zu Berlin\\
		\vspace{1.5cm}
		Robert Bärhold \& Arne Binder \\
		31. März 2014 \\ %TODO
		\vspace{1cm}
		
		\begin{table}[h]
			\Large
			\centering
			\begin{tabular}{l l}
				Seminar: & Computergestützte Analyse von Sprache\\
				Seminarleiter: & Prof. Dr. Ulf Leser\\
				 		    & Prof. Dr. Anke Lüdeling \\
				Semester: & Wintersemester 2013/14 				 	
			\end{tabular}
		\end{table}	
	\end{Large}
	\vspace*{\fill}
\end{center}


\pagenumbering{roman}
\pagestyle{fancy} %eigener Seitenstil
\fancyhf{} %alle Kopf- und Fußzeilenfelder bereinigen
\renewcommand{\headrulewidth}{0pt} %obere Trennlinie
\renewcommand{\footrulewidth}{0pt} %untere Trennlinie 
\fancyfoot[C]{\thepage} %Seitennummer

% \newpage
% \tableofcontents
% \vspace{1cm}
% \listoffigures
% \vspace{1cm}
% \listoftables

\newpage
\pagenumbering{arabic}

\section{Einleitung}
\section{Problemstellung}
\subsection{Korpus: Salsa 2.0}
\section{Realisierung der automatischen Rollenannotation}
\section{Evaluation des xxxxx} %TODO

Der entwickelte Ansatz der automatischen Rollenannotation wurde nach Abschluss
der Entwickelung einer Evaluation unterzogen. Diese sollte Aufschluss über die
Qualität der entwickelten Anwendung geben, so dass ein Vergleich mit ähnlichen
Systemen gezogen werden kann.

\newpage
\bibliography{biblio}
\bibliographystyle{natdin}


\end{document}